\chapter{Conclusion générale}
\label{chap:conclusion_generale}

L'essor récent de l'\ac{IA} et des modèles neuronaux de traduction a considérablement amélioré la qualité des systèmes de traduction automatique pour les langues disposant d'importantes ressources linguistiques. Toutefois, ces avancées profitent encore peu aux \ac{LFR}, notamment celles de la RDC, dont la faible disponibilité de corpus numériques constitue un obstacle majeur au développement d'outils de \ac{TAL}. Parmi celles-ci figure le \textbf{ruund}, qui demeure très peu représenté dans les ressources et les modèles de traduction existants. C'est dans ce contexte que s'inscrit le présent mémoire, dont l'objectif principal était d'étudier les approches de traduction automatique adaptées aux \ac{LFR} et de développer un système de traduction neuronale pour le couple \textbf{ruund--français}.

Pour répondre à cette problématique, cette recherche s'est appuyée sur une démarche méthodologique couvrant l'ensemble du cycle de développement d'un système de traduction automatique, depuis la constitution des données jusqu'au déploiement d'une application fonctionnelle. La première contribution de ce travail réside dans l'étude des fondements théoriques de la \ac{TAN} et dans l'analyse critique des recherches consacrées aux \ac{LFR}. Cette revue de la littérature a mis en évidence que, malgré les progrès réalisés grâce aux architectures Transformer et aux modèles multilingues pré-entraînés, les langues africaines restent largement sous-représentées dans les ressources linguistiques et les systèmes de traduction disponibles. Ce constat a confirmé la pertinence scientifique du choix du ruund comme langue d'étude.

La deuxième contribution majeure concerne la construction du corpus parallèle \textbf{OpenRuund}, constitué de 8~431 paires de phrases en ruund et en français. Ce corpus a été élaboré à partir de textes issus principalement du Nouveau Testament ainsi que de textes poétiques. Sa constitution a nécessité un processus rigoureux comprenant la numérisation des documents, l'extraction des textes par reconnaissance optique de caractères, le nettoyage, la normalisation, l'alignement automatique des phrases, la correction manuelle des erreurs ainsi que la validation finale des données. Cette ressource constitue, à notre connaissance, l'un des premiers corpus parallèles structurés disponibles pour le couple ruund--français et représente une contribution importante au développement futur des technologies de \ac{TAL} du ruund.

À partir de ce corpus, trois modèles de \ac{TAN}, à savoir \ac{mBART}-50, \ac{MarianMT} et \ac{NLLB}-200-distilled-600M, ont été adaptés au contexte du ruund grâce à une phase de \textit{fine-tuning}. L'absence du ruund parmi les langues initialement prises en charge par certains modèles a conduit à adapter leur vocabulaire par l'ajout du jeton linguistique \texttt{ruu\_CM} et au redimensionnement des couches d'embedding lorsque cela était nécessaire. Les performances des modèles ont ensuite été évaluées à l'aide des métriques automatiques \ac{BLEU} et \ac{chrF}, complétées par une évaluation humaine réalisée par des locuteurs natifs du ruund et un spécialiste en linguistique et littératures africaines. Les résultats obtenus montrent que les trois modèles sont capables de produire des traductions de qualité satisfaisante malgré le volume limité des données d'entraînement. Parmi eux, \ac{NLLB}-200-distilled-600M s'est distingué en obtenant les meilleures performances globales, démontrant ainsi l'intérêt des modèles multilingues massivement pré-entraînés pour les \ac{LFR}.

Au-delà de l'expérimentation des modèles, cette recherche a abouti au développement d'une application web intégrant le meilleur modèle de traduction ainsi qu'un module de collecte collaborative de données textuelles et vocales. Les modèles entraînés et le corpus construit ont également été publiés sur la plateforme Hugging Face afin de favoriser leur réutilisation par la communauté scientifique. L'ensemble de ces réalisations constitue une contribution concrète au développement de ressources numériques ouvertes pour le ruund et jette les bases de futurs travaux dans le domaine du \ac{TAL} des langues congolaises.

Ces résultats permettent de valider les hypothèses formulées au début de cette étude : la disponibilité d'un corpus parallèle de qualité constitue un facteur déterminant pour l'apprentissage de modèles de traduction performants ; l'adaptation de modèles multilingues pré-entraînés permet de compenser, dans une certaine mesure, le manque de ressources propres au ruund ; et il est possible de développer un système de traduction opérationnel pour une \ac{LFR} à partir d'un corpus de taille relativement modeste, à condition de mettre en œuvre un processus rigoureux de préparation des données et de \textit{fine-tuning}. Ces résultats confirment également que l'ensemble des objectifs fixés au début de la recherche --- étude des approches adaptées aux \ac{LFR}, constitution du corpus, adaptation et comparaison des modèles, évaluation des traductions et déploiement d'un système applicatif --- ont été atteints.

Malgré ces résultats encourageants, plusieurs difficultés ont marqué la réalisation de ce travail. La principale réside dans l'absence quasi totale de ressources numériques existantes pour le ruund. Cette situation a nécessité la constitution d'un corpus parallèle à partir de documents imprimés, impliquant un travail conséquent de numérisation et de correction des erreurs générées par la reconnaissance optique de caractères. L'alignement des phrases entre les versions ruund et française a également demandé de nombreuses vérifications manuelles afin de garantir la qualité des données utilisées pour l'entraînement. En outre, l'absence du ruund parmi les langues officiellement supportées par certains modèles pré-entraînés a nécessité l'ajout d'un nouveau jeton linguistique et l'adaptation du vocabulaire des modèles. Enfin, les contraintes liées aux ressources matérielles disponibles, notamment la capacité limitée des \ac{GPU}, ont imposé certains compromis dans les paramètres d'entraînement.

Cette étude présente également quelques limites. Le corpus parallèle utilisé, bien qu'il constitue une ressource précieuse, demeure relativement restreint comparativement aux jeux de données disponibles pour les langues à fortes ressources. De plus, une grande partie des données provient du domaine religieux, ce qui peut limiter la capacité de généralisation des modèles vers d'autres domaines tels que l'administration, l'éducation, la santé ou les médias. Par ailleurs, seuls trois modèles de \ac{TAN} ont été comparés dans cette étude, tandis que d'autres architectures récentes auraient pu être évaluées. Enfin, l'évaluation humaine repose sur un nombre limité d'évaluateurs, ce qui ne permet pas d'apprécier de manière exhaustive la qualité des traductions produites.

Ces limites ouvrent plusieurs perspectives de recherche. Une première perspective consiste à enrichir progressivement le corpus OpenRuund en intégrant des textes provenant de domaines variés afin d'améliorer la couverture lexicale et la capacité de généralisation des modèles. La poursuite de la collecte collaborative de données textuelles et vocales, à travers la plateforme développée dans ce travail, permettra également d'augmenter le volume des ressources disponibles tout en impliquant davantage de locuteurs natifs. Sur le plan méthodologique, il serait pertinent d'explorer des techniques récentes d'adaptation efficace des modèles, telles que PEFT, LoRA ou QLoRA, ainsi que l'utilisation de modèles multilingues de nouvelle génération. Une autre perspective importante consisterait à étendre cette approche à d'autres langues congolaises à faibles ressources afin de favoriser le développement d'un écosystème national de technologies linguistiques. Enfin, le système développé pourrait évoluer vers une plateforme de traduction multimodale intégrant la reconnaissance automatique de la parole et la synthèse vocale, ouvrant ainsi la voie à des applications de traduction parole-à-parole.

En définitive, ce mémoire démontre que le développement d'un système de traduction automatique pour une \ac{LFR} telle que le ruund est non seulement possible, mais également prometteur lorsque celui-ci s'appuie sur une démarche méthodique de constitution des données, d'adaptation de modèles multilingues et d'évaluation rigoureuse des performances. Au-delà des résultats scientifiques obtenus, cette recherche contribue à l'enrichissement des ressources numériques disponibles pour le ruund, à la valorisation du patrimoine linguistique congolais et à la promotion d'une \ac{IA} plus inclusive. Les ressources produites, les modèles entraînés et la plateforme développée constituent une base solide sur laquelle pourront s'appuyer de futurs travaux visant à accélérer l'intégration des langues congolaises dans les technologies modernes de \ac{TAL}.